\documentclass[10pt,a4paper]{amsart}
\usepackage[margin=19mm]{geometry}
\usepackage{amsmath,amssymb,mathtools}
\usepackage[scr=rsfs,cal=euler]{mathalfa}
\usepackage{xfrac}
\usepackage[unicode,hidelinks,bookmarks=false]{hyperref}
\newcommand{\ie}{i.e.\ }
\newcommand{\cf}{cf.\ }
\newcommand{\resp}{respectively\ }
\newcommand{\Subsection}{\subsection}
\providecommand{\nobreakdash}{\nobreak\hbox{-}}
\setlength{\emergencystretch}{2em}
\multlinegap0pt
\usepackage{amssymb}
\newtheorem{lem}{}
\title{On the II$_{1}$ Factors of Fuchsian Groups}
\author{D. Shlyakhtenko}
\date{Source: arXiv:2609.11074v1 (CC BY 4.0)}
\begin{document}
\maketitle

\noindent\emph{Source and licence.} On the II$_{1}$ Factors of Fuchsian Groups. Version arXiv:2609.11074v1 (CC BY 4.0), distributed by arXiv under the Creative Commons Attribution 4.0 International licence, http://creativecommons.org/licenses/by/4.0/, as recorded in the arXiv licence metadata for this exact version and shown on the abstract page https://arxiv.org/abs/2609.11074v1. Full licence notice: \texttt{licenses/arxiv-2609.11074-CC-BY-4.0.txt}.\par\smallskip

\noindent\textit{Editorial scope.} Counted unit: the article's lem lem:F\_tlowerEqThetaOfB (source0.tex lines 970--976). The article's complete original proof of it is delivered with it (source0.tex lines 977--1170, one \texttt{proof} environment of 194 lines). No mathematics is written, repaired or re-derived: the statement and the proof are the article's own lines, in the article's own notation, and no retained line is substituted or re-flowed.  Retained uncounted as the source-local context the proof needs: lines 506--514; lines 946--953.  Nothing was omitted from any retained span, so the package carries no omission record.  No retained line contains a citation, so the wrapper carries no bibliography and no bibliography entry.  No retained line contains a figure environment, an image-inclusion command or drawing code, so no image is delivered and the package declares no asset dependency.  Editorial formatting only, in addition: an AMS \texttt{amsart} preamble that restates the article's own declarations and macros used by the retained lines, namely the environments \texttt{lem} on separate counters, together with the standard packages those lines need.  The retained environments print in this document as Lemma 1 in order of appearance, whereas the article prints the same items with the numbering of its own full document.  The complete native source of the article as distributed in the arXiv e-print for this exact version is bundled unchanged as \texttt{sources/209986/source0.tex}.
\begin{lem}
\label{lem:descriptionOfSfWords}Let $g\in\mathbb{F}_{2}$ be a reduced
word containing $p\geq1$ occurrences of $A^{\pm1}$, $f:\mathbb{\mathbb{T}\to\mathbb{T}}$
measurable, and let $c_{j}(k)$, $W$, and $K$ be as above. Then:
\begin{itemize}
\item If $h$ is not in the range of the map $W$, then $\langle h,S_{f}(g)\rangle=0.$
\item If $h=W(k_{0},\dots,k_{p})$, then $\langle h,S_{f}(g)\rangle=\prod_{j=0}^{p}c_{j}(k_{j}).$
\end{itemize}
\end{lem}

\begin{lem}
\label{lem:propertiesOffmt}(i) $f_{m,t}:\mathbb{T}\to\mathbb{T}$
is measurable. (ii) Let $\ensuremath{c_{m}(k)=\int_{\mathbb{T}}\bar{f}_{m,t}(z)z^{-k}d\mu(z)=\widehat{\bar{f}_{m,t}}(k)}$,
and let $\alpha_{m}=\left|\widehat{\bar{f}_{m,t}}(0)\right|^{2}$,
$\beta_{m}=\left|\widehat{\bar{f}_{m,t}}(\pm1)\right|^{2}$. Then
$\alpha_{m}=1-(t/m)+O(m^{-2})$ and $\beta_{m}=(t/2m)+O(m^{-2})$.
(iii) The Fourier series for $f$ is absolutely convergent. 
\end{lem}

\begin{lem}
\label{lem:F_tlowerEqThetaOfB}For every $t>0$ and every $\epsilon>0$
there exists an $m_{0}$ so that for all $m>m_{0}$ and all $0\leq q\leq1$,
\[
Q_{q}(\theta_{m,t}^{-1}(B))\geq F_{t}(q)-\epsilon.
\]
\end{lem}

\begin{proof}
Let $\ensuremath{c_{m}(k)=\int_{\mathbb{T}}\bar{f}_{m,t}(z)z^{-k}d\mu(z)=\widehat{\bar{f}_{m,t}}(k)}$,
and let $\alpha_{m}=\left|\widehat{\bar{f}_{m,t}}(0)\right|^{2}$,
$\beta_{m}=\left|\widehat{\bar{f}_{m,t}}(\pm1)\right|^{2}$ as in
Lemma \ref{lem:propertiesOffmt}; thus $\alpha_{m}=1-(t/m)+O(m^{-2})$
and $\beta_{m}=(t/2m)+O(m^{-2})$. 

Since $\theta_{m,t}^{-1}=R_{-m}\circ S_{\bar{f}_{m,t}}\circ R_{m}$,
it follows from commutation between $T_{q}$ and $R_{m}$ that $Q_{q}(\theta_{m,t}^{-1}(B))=Q_{q}(S_{\bar{f}_{m,t}}(R_{m}(B)))=Q_{q}(S_{\bar{f}_{m,t}}(BA^{m}))$.

Lemma \ref{lem:descriptionOfSfWords} then gives us
\[
S_{\bar{f}_{m,t}}(BA^{m})=\sum_{\mathbf{n}\in\mathbb{Z}^{m}}\left(\prod_{j=1}^{m}c_{m}(n_{j})\right)W(\mathbf{n})
\]
where $\mathbf{n}\to W(\mathbf{n})$ is the injective map
\[
W(\mathbf{n})=BAB^{n_{1}}AB^{n_{2}}\cdots AB^{n_{m}}.
\]
Thus
\begin{align}
Q_{q}(\theta_{m,t}^{-1}(B)) & =\sum_{\mathbf{n}\in\mathbb{Z}^{m}}\left(\prod_{j=1}^{m}|c_{m}(n_{j})|^{2}\right)q^{\ell(W(\mathbf{n}))}\nonumber \\
 & \geq\sum_{\mathbf{n}\in\{-1,0,1\}^{m-1}\times\{0\}}\left(\prod_{j=1}^{m}|c_{m}(n_{j})|^{2}\right)q^{\ell(W(\mathbf{n}))}.\label{eq:termsWeKeep}
\end{align}
We will use the smaller sum (keeping only terms with $n_{1},\dots,n_{m-1}\in\{-1,0,1\}$
and $n_{m}=0$) for our lower estimate. Thus from now on we only consider
$n_{m}=0$ and each $n_{1},\dots,n_{m-1}\in\{-1,0,1\}$. Our aim is
to bound $\ell(W(\mathbf{n})).$

Recall that $H$ is generated by $A$ and $C=BAB^{-1}$. 

Define $P_{0}=BA$, $P_{j}=P_{j-1}B^{n_{j}}A$, $1\leq j\leq m-1$.
Since $n_{m}=0$,
\[
W(\mathbf{n})=P_{m-1}.
\]
We recursively factor
\[
P_{j}=V_{j}B^{e_{j}},\qquad e_{j}\in\{0,1\}
\]
as follows. Initially, $P_{0}=BA=CB$, thus we set $V_{0}=C$ and
$e_{0}=1$. Suppose we factored $P_{j}=V_{j}B^{e_{j}}$. Then
\[
P_{j+1}=P_{j}B^{n_{j+1}}A.
\]
We then set $V_{j+1},e_{j+1}$ according to the table:
\[
(V_{j+1},e_{j+1})=\begin{cases}
(V_{j}B^{-1}A,0), & e_{j}+n_{j+1}=-1,\\
(V_{j}A,0) & e_{j}+n_{j+1}=0,\\
(V_{j}C,1) & e_{j}+n_{j+1}=1,\\
(V_{j}BC,1) & e_{j}+n_{j+1}=2.
\end{cases}
\]
Because $BA=CB$ and $B^{2}A=BCB$, $P_{j+1}=V_{j+1}B^{e_{j+1}}$.
We claim that then $P_{j+1}=V_{j+1}B^{e_{j+1}}$. Indeed, there are
6 possible steps:
\begin{itemize}
\item (Type 1) Suppose that $e_{j}+n_{j+1}=-1$ , so that $e_{j}=0$ and
$n_{j+1}=-1$. Then $P_{j}=V_{j}B^{0}$ , $V_{j+1}=V_{j}B^{-1}A$,
and $P_{j+1}=P_{j}B^{-1}A=V_{j+1}B^{0}$.
\item Suppose that $e_{j}+n_{j+1}=0$. There are two possibilities:
\begin{itemize}
\item (Type 2) If $e_{j}=0$ and $n_{j+1}=0$ , then $P_{j}=V_{j}B^{o}$,
$V_{j+1}=V_{j}A$, and $P_{j+1}=P_{j}A=V_{j+1}B^{0}.$
\item (Type 3) If $e_{j}=1$ and $n_{j+1}=-1$, then $P_{j}=V_{j}B$, $V_{j+1}=V_{j}A$,
and $P_{j+1}=P_{j}B^{-1}A=V_{j}A=V_{j+1}B^{0}.$
\end{itemize}
\item Suppose that $e_{j}+n_{j+1}=1$. There are two possibilities:
\begin{itemize}
\item (Type 4) If $e_{j}=0$ and $n_{j+1}=1$, then $P_{j}=V_{j}B^{0}$,
$V_{j+1}=V_{j}C=V_{j}BAB^{-1}$, and $P_{j+1}=P_{j}BA=V_{j+1}B.$
\item (Type 5) If $e_{j}=1$ and $n_{j+1}=0$, then $P_{j}=V_{j}B$, $V_{j+1}=V_{j}C=V_{j}BAB^{-1}$,
and $P_{j+1}=P_{j}A=V_{j+1}B.$
\end{itemize}
\item (Type 6) Suppose finally that $e_{j}+n_{j+1}=2$ , so that $e_{j}=1$
and $n_{j+1}=1$ . Then $P_{j}=V_{j}B$, $V_{j+1}=V_{j}BC=V_{j}B^{2}AB^{-1},$and
$P_{j+1}=P_{j}BA=V_{j+1}B.$
\end{itemize}
Let us say that the factorization $P_{j}=V_{j}B^{e_{j}}$ is even
(respectively, odd) if $e_{j}=0$ (resp., $e_{j}=1$).

If the transition from $P_{j}=V_{j}B^{e_{j}}$ to $P_{j+1}=V_{j+1}B^{e_{j+1}}$
is of type 2 or 5, then the parity is unchanged, and $\ell(V_{j+1})=\ell(V_{j})$
since $V_{j}$ is multiplied by an element of $H$. We call these
transitions \textbf{type A}. In this case, $n_{j+1}=0$.

If the transition is of type 3 or 4, then the parity is reversed,
but $\ell(V_{j+1})=\ell(V_{j})$ since $V_{j}$ is multiplied by an
element of $H$. We call these transitions \textbf{type B}. In this
case, $|n_{j+1}|=1$.

Finally, if the transition is of type 1 or 6, then the parity is unchanged,
but $V_{j}$ is changed by the multiplication with either $B$ or
$B^{-1}$ as well as an element of $H$. Thus the length of $V_{j+1}$
could have increased by $1$. We call these transitions \textbf{type
C}. In this case, $|n_{j+1}|=1$.

Let $a_{j}$ , $b_{j}$, and $c_{j}$ be the total number of type
A, type B, and type C transitions that occurred when moving from $P_{0}=V_{0}B$
to $P_{j}=V_{j}B^{e_{j}}$ . Then
\[
\ell(V_{j})\leq c_{j},\qquad\ell(P_{j})\leq c_{j}+e_{j}.
\]
Let $a=a_{m-1}$, $b=b_{m-1}$, $c=c_{m-1}$ be the total number of
transitions of each type. Then $a+b+c=m-1$. Since only type B transitions
affect parity, and the inital parity is odd, the final parity is odd
iff $b$ is even. 
\[
\ell(W(\mathbf{n}))=\ell(P_{m-1})\leq c+1_{b\textrm{ even}}.
\]
Moreover, the number of indices $j$ for which $|n_{j}|=1$ is precisely
$c+b$ and the number of $j$ for which $n_{j}=0$ is precisely $a$.
Therefore, the term in (\ref{eq:termsWeKeep}) associated to this
particular $\mathbf{n}$ is 
\begin{equation}
\alpha_{m}^{a+1}\beta_{m}^{b+c}q^{\ell(W(\mathbf{n}))}\geq\alpha_{m}^{a+1}\beta_{m}^{b+c}q^{c+1_{b\textrm{ even}}}.\label{eq:estimateOnabTerm}
\end{equation}

Note also that if we are given the sequence of transition types, then
we can recover $\mathbf{n}$ fully. Indeed, the value of $n_{j+1}$
is completely determined if we know the transition type and the parity
$e_{j}$: for transitions of type A, $n_{j+1}=0$; for type B, $n_{j+1}$
is negative iff $e_{j}=1$, and for type C, $n_{j+1}$ is positive
iff $e_{j}=1$. Since $e_{0}=1$, it follows that the sequence of
transition types completely determines $\mathbf{n}$. Conversely,
knowing $\mathbf{n}$ produces a sequence of $m-1$ transition types.
Therefore, there is a bijection between sequences of $m-1$ transition
types and tuples $\mathbf{n}$. One can realize this transition explicitly
using a generating function.

Consider the polynomial $\mathscr{P}_{m}(X,Y)$ given by
\[
\mathscr{P}_{m}(X,Y)=(\alpha_{m}+\beta_{m}X+\beta_{m}Y)^{m-1}\alpha_{m}.
\]
If one expands the right hand side of this equation, then each term
is a product of $m$ terms, each chosen from among $\alpha_{m}$,
$\beta_{m}X$ and $\beta_{m}Y$. Keeping the order of terms in this
product, we see that the products are in bijection with sequences
of $m-1$ of three types, with the three terms corresponding to types
A, B, and C. Moreover, $b$ and $c$ correspond precisely to the total
power of $X$ and $Y$ in the corresponding term.

Let $p_{b,c}^{(m)}$ be the coefficient of $X^{b}Y^{c}$ in $\mathscr{P}_{m}(X,Y)$.
Then
\[
p_{b,c}^{(m)}=\sum_{\mathbf{n}\textrm{ has \ensuremath{b} type B and \ensuremath{c} type C}}\alpha_{m}^{a+1}\beta_{m}^{b+c},
\]
where $a=m-b-c-1.$ Therefore, (\ref{eq:termsWeKeep}) combined with
(\ref{eq:estimateOnabTerm}) implies that
\[
Q_{q}(\theta_{m,t}^{-1}(B))\geq q\sum_{b,c\geq0\textrm{ with \ensuremath{b} even}}p_{b,c}^{(m)}q^{c}+\sum_{b,c\geq0\textrm{ with \ensuremath{b} odd}}p_{b,c}^{(m)}q^{c}.
\]
These sums can be recovered by considering the even and odd parts
of $\mathscr{P}$ and substituting $X=q$ and $Y=\pm1$: 
\[
\sum_{b,c\textrm{ with \ensuremath{b} even}}p_{b,c}^{(m)}q^{c}=\frac{1}{2}(\mathscr{P}_{m}(1,q)+\mathscr{P}_{m}(-1,q))
\]
and
\[
\sum_{b,c\textrm{ with \ensuremath{b} odd}}p_{b,c}^{(m)}q^{c}=\frac{1}{2}(\mathscr{P}_{m}(1,q)-\mathscr{P}_{m}(-1.-q)).
\]
Substituting this in gives us the final estimate
\[
Q_{q}(\theta_{m,t}^{-1}(B))\geq\frac{q+1}{2}\mathscr{P}_{m}(1,q)+\frac{q-1}{2}\mathscr{P}_{m}(-1,q).
\]

Let $\sigma\in\{-1,1\}$. We now substitute $\alpha_{m}=1-(t/m)+O(m^{-2})$
and $\beta_{m}=(t/2m)+O(m^{-2})$ into the equation for $\mathscr{P}_{m}$:
\[
\mathscr{P}_{m}(q,\sigma)=(\alpha_{m}+\beta_{m}q+\beta_{m}\sigma)^{m-1}\alpha_{m}.
\]
For $m$ large enough $\alpha_{m}+\beta_{m}(q+\sigma)=1-(t/m)+(t/2m)(q+\sigma)+O(m^{-2})$
(where $O(m^{-2})$ is uniform in $q$). In particular, for large
enough $m$ this is positive for all $q\in[-1,1]$ and both choices
of $\sigma$. Taking logarithms of both sides gives us
\begin{align*}
\log\mathscr{P}_{m}(q,\sigma) & =(m-1)\log(1+\frac{t}{2m}(q+\sigma-2)+O(m^{-2}))+\log(1-(t/m)+O(m^{-2}))\\
 & =(m-1)\frac{t}{2m}(q+\sigma-2)+O(m^{-1})
\end{align*}
uniformly in $q$ (and pointwise in $t$). Thus
\[
\mathscr{P}_{m}(q,\sigma)\to\exp\left(\frac{t}{2}(q+\sigma-2)\right).
\]
It follows that
\[
\frac{q+1}{2}\mathscr{P}_{m}(q,1)+\frac{q-1}{2}\mathscr{P}_{m}(q,-1)\to e^{t(q-1)/2}\left(\frac{1-e^{-t}}{2}+\frac{1+e^{-t}}{2}q\right)=F_{t}(q).
\]
Thus given $\epsilon>0$ we can choose $m_{0}$ so that for any $m>m_{0}$,
for all $q\in[0,1]$, 
\[
Q_{q}(\theta_{m,t}^{-1}(B))\geq F_{t}(q)-\epsilon,
\]
as claimed.
\end{proof}

\end{document}
